\chapter{Мешок с деталями}
% полная версия: chapters/01b_neurontypes.tex

\pencilfig{1}{\pencilart{1}}{Мешок с деталями: почти все синие, немного красных и горстка остальных.}

Представьте, что вам дали мешок, в котором лежат восемьдесят шесть миллиардов одинаковых на вид деталей, и попросили разобраться, как из них собран компьютер. С чего бы вы начали? Инженер не стал бы разглядывать форму деталей. Он спросил бы: сколько у детали входов, сколько выходов и что она выдаёт на выход.

Нейрон --- так называется такая деталь --- устроен удивительно просто. Входов у него тысячи: к нему подходят провода от других нейронов. Выход один. Но этот выход ветвится, как дерево, и одна и та же копия сигнала уходит к тысячам получателей. Представьте человека, который слушает тысячу собеседников сразу, а отвечает одной фразой, но так, что её слышат ещё тысяча человек.

Что это за сигнал? Короткий щелчок. Нейрон время от времени выдаёт импульс --- всплеск напряжения длиной около тысячной доли секунды. Все щелчки одного нейрона одинаковы, как удары одного барабана. Поэтому всё, что нейрон может сказать, записано не в громкости щелчка, а в том, \emph{когда} он щёлкает: чаще, реже, сразу или с задержкой.

Внутри нейрон делает самое простое вычисление, какое можно придумать: складывает приходящие сигналы и, если сумма перевалила за порог, щёлкает сам. Одни входы толкают сумму вверх, другие --- вниз.

\section*{Как рассортировать мешок}

Кончик выходного провода, дотянувшись до соседа, не касается его. Между ними остаётся узкая щель, и сигнал через неё передаёт химия: провод выбрасывает щепотку особого вещества --- медиатора, --- а сосед его ловит. И вот удача: почти каждый нейрон на всех кончиках своего провода выбрасывает одно и то же вещество. Значит, нейроны можно сортировать по тому, \emph{что} они выбрасывают.

Мы будем называть тип нейрона первыми четырьмя буквами его вещества. Нейрон, выбрасывающий глутамат, --- \nt{ГЛУТ}. Выбрасывающий гамма-аминомасляную кислоту --- \nt{ГАМК}. Дальше пойдут \nt{ДОФА} (дофамин), \nt{СЕРО} (серотонин), \nt{НОРА} (норадреналин), \nt{АЦЕТ} (ацетилхолин) и ещё несколько.

Когда мешок рассортирован, картина получается очень неравной. Из каждой тысячи нейронов около девятисот шестидесяти --- это \nt{ГЛУТ}. Их сигнал подталкивает получателя к щелчку: это «плюс». Ещё около тридцати шести --- \nt{ГАМК}, их сигнал отталкивает от порога: это «минус». В коре их больше, от пятой части до четверти. А все остальные типы вместе --- несколько нейронов на сто тысяч. Капля в море.

\section*{Телефон и радио}

Но эта капля устроена совсем иначе. Провод \nt{ГЛУТ}- или \nt{ГАМК}-нейрона ведёт к определённым соседям, как телефонная линия: сигнал получают только те, кому звонят. А крошечные группы \nt{ДОФА}, \nt{СЕРО}, \nt{НОРА} и \nt{АЦЕТ} работают как радиостанции. Их провода расходятся по огромным областям мозга, вещество часто выливается не в щель, а в пространство вокруг, и одно и то же сообщение слышат сразу миллионы клеток.

По телефонным линиям идут данные: вот линия на картинке, вот звук определённой высоты. По радио идут общие настройки: насколько всё громко, стоит ли учиться, пора ли спать. В любом компьютере так же: ячеек памяти миллиарды, а управляющих регистров --- горстка. Без них машина всё равно не работает.

\section*{Смысл выбирает получатель}

Есть одна тонкость. Само вещество не знает, «плюс» оно или «минус». Молекула --- это ключ. Что произойдёт, когда ключ вставят, решает замок --- особый белок на стороне получателя. Один и тот же дофамин одних получателей подбадривает, а других приглушает, смотря какой замок у них стоит. Это как одна и та же радиопередача, которую разные слушатели понимают по-разному.

\section*{Сигнал «лучше, чем ожидалось»}

Самый знаменитый радиоканал --- дофамин. Что он передаёт? Вот опыт. Животному неожиданно дают каплю сока, и дофаминовые нейроны отвечают вспышкой щелчков. Потом перед соком каждый раз зажигают лампочку. Животное выучивает, что лампочка обещает сок, и теперь вспышка случается на лампочку, а на сам сок --- нет. А если после лампочки сок не дают, дофаминовые нейроны в момент ожидания \emph{замолкают}.

Правило, которое объясняет все три случая: дофамин сообщает не «хорошо», а «лучше, чем я ожидал». Неожиданный сок --- приятный сюрприз. Обещанный --- никакого сюрприза. Обещанный, но не данный --- неприятный. Ровно такую величину используют программы, которые учатся на своих ошибках. К ней мы ещё вернёмся.

\section*{Что унести с собой}

Мозг --- это огромная сеть одинаковых на вид деталей. Почти все они --- «плюсы» и «минусы» телефонной сети данных. И совсем немногие --- радиостанции, которые настраивают всю сеть сразу. В следующей главе мы проверим, что можно построить из одних только «плюсов».

\fullversion{1}
